% BEGIN GENERATED ARTIFACT: prf:existence-disjunctive-normal-form-propositional-logic
\newpage
\phantomsection
\label{prf:existence-disjunctive-normal-form-propositional-logic}
\LRAProofFor{thm:existence-disjunctive-normal-form-propositional-logic}

\begin{remark*}[Return]
\hyperref[thm:existence-disjunctive-normal-form-propositional-logic]{Return to Theorem (Existence of Disjunctive Normal Form).}
\end{remark*}

\begin{theorem*}[Existence of Disjunctive Normal Form]
For every formula \(\varphi\in\WFF\), there exists a formula
\(\psi\in\WFF\) such that \(\psi\) is in disjunctive normal form and
\[
\varphi\equiv\psi.
\]
\end{theorem*}

\begin{proof}
\textbf{Professional Standard Proof.}~
TODO: supply the compact proof.
\end{proof}

\begin{proof}
\textbf{Detailed Learning Proof.}~
TODO: supply the detailed learning proof.
\end{proof}

\begin{remark*}[Proof structure]
First use the existence theorem for negation normal form to replace
\(\varphi\) by an equivalent NNF formula. Then prove, by induction on that NNF
formula, that repeated use of the distributive laws transforms it into a
disjunction of elementary conjunctions. Replacement of logical equivalents
preserves equivalence through each local transformation.
\end{remark*}

\begin{dependencies}
\begin{itemize}
  \item \hyperref[def:disjunctive-normal-form-propositional-logic]{Disjunctive Normal Form}
  \item \hyperref[def:negation-normal-form-propositional-logic]{Negation Normal Form}
  \item \hyperref[thm:distributive-laws-propositional-logic]{Distributive Laws for Propositional Logic}
  \item \hyperref[thm:replacement-of-logical-equivalents-propositional-logic]{Replacement of Logical Equivalents}
  \item \hyperref[thm:existence-negation-normal-form-propositional-logic]{Existence of Negation Normal Form}
\end{itemize}
\end{dependencies}

\clearpage
% END GENERATED ARTIFACT: prf:existence-disjunctive-normal-form-propositional-logic
